% !TeX root = ../papers/nc-p3.tex
\section{固定曲線の倍率を動かしたときの長さ5}
\label{sec:varying-fixed-quintics}
\begin{theorem}\label{thm:varying-fixed-quintics}
負の固定曲線では全通常倍率で、正の固定曲線では $(3,2)$ を含む倍率平面の
非空Zariski開集合で、長さ4を全始点で通過した全ての許容窓の長さ5の次元は全始点で56である。
\end{theorem}
\begin{proof}
同時文字交換の自己同型は各矢の次数を $j\mapsto j\mathbin\oplus1$ に送る。
従って奇数長の道で次数 $h$ と $h\mathbin\oplus1$ の次元が等しい。
二位置移動も併用すると、正負ごとに始点0,1、文字0,2を検算すればよい。
前節の行列から各一つの階数86の小行列式を取り、
$e=Nu,f=Du$ を入れて通常倍率で単元の分母を払う。
\path{computations/verify_varying_fixed_quintics.py} はこの行列式を
$\mathbb Z[a,u,z,\lambda]$ で因数分解する。具体的には非零成分を掃き出す際、
ピボット $p$ と消去成分 $x$ の最大公約数を $g$ とし、
行を $(p/g)R-(x/g)R_{\rm pivot}$ に置き換える。
各行の内容を割り、取り除いた因子と $p/g$ の逆倍率を記録する。
行・列の交換の符号と有理定数を除いた行列式の全因子を復元できる。

各因子を $G$ との $u$ に関する終結式で第一座標へ写す。
負では四つの検算の全てで、非定数因子は定理\ref{thm:uniform-fixed-quartic}
の表と通常倍率で逆にした因子に限られる。全根が四次で失格する座標と
非許容端点を除けば、従って全通常倍率で目標階数に達する。

正の側では、同定理の第一座標の表以外にも因子 $P(a,z,\lambda)$ が出る。
これらについて検算は個々に次の三条件を確認する。
\begin{enumerate}
\item 倍率だけに依存する追加因子はない。
\item $(3,2)$への特殊化は $a$ の次数を保ち、平方自由である。
\item 特殊化の全根で、$a\Delta ND$ と $u$ の二次式 $G$ の最高次係数は非零である。
既知の四次障害の第一座標はなく、$G$ と $xy\det B$ の終結式も消えない。
\end{enumerate}
これにより $(3,2)$ の近傍で $P=0$ は倍率平面に有限で、
その上の曲線 $G=0$ も有限になる。必要に応じて基底を縮め、上記の非零条件を保つ。
この有限族の中で目標階数に達しない閉集合の像は閉じる。
命題\ref{prop:fixed-quintic-slice}によりその像は $(3,2)$ を含まない。

表の中の第一座標についても注意する。正の古い二座標の根は四次で失格する。
座標端とその交換相手には、許容性の下で一意な有理表示がある。
二つの $H_\epsilon$ の根では、反対行列式の既知の点の相方を
定理\ref{thm:uniform-fixed-quartic}の式で比較する。$(3,2)$で
これらの四つの第一座標は相異なり $abND\ne0$ である。
相方では $xy\det B(\det A+\det B)\ne0$ である
（定理の四つの残余表示に $r=\pm\sqrt3,\lambda=2$ を代入すればよい）。
従って全ての相方と許容端点で、前節の長さ5の階数が開近傍に延びる。
$r^2=z$の有限被覆から悪い閉集合の像を除けば、この近傍は倍率平面上に取れる。
こうして残る全ての第一座標を同時に扱える。
\end{proof}
